\section{Introducción}
La inferencia de redes neuronales contiene operaciones que multiplican
activaciones por pesos y reducen sus productos. Cuantizar estos tensores
permite representar el cálculo mediante enteros y parámetros de escala,
pero su utilidad depende tanto del formato como de su implementación~\cite{jacob}.
En un core escalar, una palabra de 32 bits puede almacenar ocho pesos de
cuatro bits; consumirlos puede exigir máscaras, desplazamientos, corrección
de origen y multiplicaciones.

Una extensión que realice parte de ese trabajo puede reducir instrucciones.
Sin embargo, añade lógica y condiciona el control del pipeline. La reducción
del conteo de instrucciones no garantiza menor tiempo si crecen la latencia
o el período de reloj. También existe una alternativa algebraica: corregir
el producto punto completo usando una suma de activaciones reutilizable.
Esta alternativa impide atribuir automáticamente un beneficio a la fusión
del zero-point.

Se propone estudiar XQDot4Z sobre kirky-arqui, un procesador docente con
pipeline de cinco etapas y soporte parcial de instrucciones RV32I/C. El objeto
experimental es el producto punto y la matriz--vector cuantizados de batch
uno. La operación no pretende ejecutar por sí sola una red completa ni
resolver todos los costos de un modelo de lenguaje.

Las contribuciones candidatas son: (i) un contrato e implementación trazables;
(ii) una comparación que distinga multiplicación, procesamiento empacado y
corrección fusionada; y (iii) una caracterización de las condiciones que
favorecen o perjudican esa fusión. Esta formulación admite resultados negativos:
si la corrección factorizada resulta superior, explicar por qué responderá
parte de la pregunta de investigación.
